\documentclass[aspectratio=169,11pt]{beamer}
\usetheme{Madrid}
\usecolortheme{seahorse}
\setbeamertemplate{navigation symbols}{}
\setbeamertemplate{footline}[frame number]
\usepackage[T1]{fontenc}
\usepackage{amsmath,amssymb,mathtools}
\usepackage{listings}
\lstset{basicstyle=\ttfamily\small, columns=fullflexible, keepspaces=true,
  showstringspaces=false, frame=single, xleftmargin=2pt, xrightmargin=2pt}
\newcommand{\Bool}{\mathsf{Bool}}
\newcommand{\Int}{\mathsf{Int}}
\newcommand{\Arrow}{\mathbin{\to}}
\newcommand{\judge}[3]{#1\vdash #2:#3}
\title[Type-checking]{Type-checking: making sure programs make sense}
\subtitle{A friendly tour of simple types}
\author{CS 413}
\date{October 20, 2026}

\begin{document}

\begin{frame}
  \titlepage
\end{frame}

\begin{frame}[fragile]{A tiny puzzle}
\begin{columns}[T]
\column{.52\textwidth}
\begin{block}{This looks fine}
\begin{lstlisting}
if 3 < 5 then 10 else 20
\end{lstlisting}
\end{block}
\begin{block}{This one?}
\begin{lstlisting}
if 3 then false else 20
\end{lstlisting}
\end{block}
\column{.43\textwidth}
Can we catch the problem before running the program?\medskip

Type-checking answers questions like:
\begin{itemize}
  \item Is the condition a boolean?
  \item Do both branches agree on a result type?
  \item Is a function given the kind of input it expects?
\end{itemize}
\end{columns}
\end{frame}

\begin{frame}{What do types buy us?}
\begin{itemize}
  \item Catch certain mistakes early, with a useful location and explanation.
  \item Document how pieces of a program are meant to fit together.
  \item Let later stages make safe assumptions (for example, about operations).
\end{itemize}
\bigskip
\begin{alertblock}{A type checker is not an oracle}
It does not prove that a program is useful, terminates, or matches the programmer's intent. It checks a particular set of consistency rules.
\end{alertblock}
\end{frame}

\begin{frame}{Our little language}
\begin{columns}[T]
\column{.49\textwidth}
\textbf{Expressions}
\begin{align*}
e ::= {}& n \mid \mathsf{true} \mid \mathsf{false} \mid x \mid e+e \mid e<e\\
       &\mid \mathsf{if}\ e\ \mathsf{then}\ e\ \mathsf{else}\ e
        \mid \lambda x{:}T.e \mid e\ e
\end{align*}
\column{.47\textwidth}
\textbf{Types}
\[
T ::= \Int \mid \Bool \mid T\Arrow T
\]
\medskip
Examples:
\begin{itemize}
  \item $42 : \Int$
  \item $\mathsf{false} : \Bool$
  \item $\Int\Arrow\Bool$ is the type of functions from integers to booleans.
\end{itemize}
\end{columns}
\end{frame}

\begin{frame}{A type is a promise about use}
\begin{block}{Read the arrow left to right}
\[
  \Int\Arrow\Bool \qquad\text{means}\qquad
  \text{“give me an integer; I'll return a boolean.”}
\]
\end{block}
\begin{exampleblock}{A typed function}
\[
  \lambda (x:\Int).\ x<10 : \Int\Arrow\Bool
\]
The annotation says what the function expects; the body confirms what it returns.
\end{exampleblock}
\begin{itemize}
  \item Function types associate to the right: $\Int\Arrow\Bool\Arrow\Int$ means $\Int\Arrow(\Bool\Arrow\Int)$.
\end{itemize}
\end{frame}

\begin{frame}{Typing judgments: the checker’s question}
\[
  \Gamma \vdash e:T
\]
Read: “under assumptions $\Gamma$, expression $e$ has type $T$.”
\begin{itemize}
  \item $\Gamma$ is the environment: a lookup table for variables.
  \item Example: $\Gamma = \{x:\Int,\ f:\Int\Arrow\Bool\}$.
  \item Then $\Gamma\vdash f\ x:\Bool$.
\end{itemize}
\begin{block}{Why keep assumptions around?}
When checking a function body, we need to remember the parameter's type. When checking a variable, we look it up.
\end{block}
\end{frame}

\begin{frame}{The easy rules}
\small
\begin{columns}[T]
\column{.49\textwidth}
\textbf{Constants and variables}
\[
\frac{}{\Gamma\vdash n:\Int}\qquad
\frac{}{\Gamma\vdash \mathsf{true}:\Bool}\qquad
\frac{x:T\in\Gamma}{\Gamma\vdash x:T}
\]
\textbf{Addition}
\[
\frac{\Gamma\vdash e_1:\Int\quad \Gamma\vdash e_2:\Int}
     {\Gamma\vdash e_1+e_2:\Int}
\]
\column{.49\textwidth}
\textbf{Comparison}
\[
\frac{\Gamma\vdash e_1:\Int\quad \Gamma\vdash e_2:\Int}
     {\Gamma\vdash e_1<e_2:\Bool}
\]
\textbf{Conditional}
\[
\frac{\Gamma\vdash c:\Bool\quad \Gamma\vdash a:T\quad \Gamma\vdash b:T}
     {\Gamma\vdash \mathsf{if}\ c\ \mathsf{then}\ a\ \mathsf{else}\ b:T}
\]
\end{columns}
\medskip
The line reads “if the things above it check, the thing below it checks.”
\end{frame}

\begin{frame}{Functions: introduce, then use}
\small
\textbf{Function introduction} (check a lambda):
\[
\frac{\Gamma,x:S\vdash e:T}
     {\Gamma\vdash \lambda(x:S).e:S\Arrow T}
\]
Assume the parameter has type $S$, then check the body.
\medskip

\textbf{Function application} (check a call):
\[
\frac{\Gamma\vdash f:S\Arrow T\quad \Gamma\vdash a:S}
     {\Gamma\vdash f\ a:T}
\]
The argument type must match the input type; the result has the output type.
\begin{exampleblock}{Quick check}
If $f:\Int\Arrow\Bool$, then $f\ 7:\Bool$ is fine, while $f\ \mathsf{false}$ is not.
\end{exampleblock}
\end{frame}

\begin{frame}{Work one through: an integer test}
Check $\lambda(x:\Int).\ x<10$.
\[
\frac{\displaystyle
  \frac{\displaystyle
    \frac{}{x:\Int\vdash x:\Int}\qquad
    \frac{}{x:\Int\vdash 10:\Int}}
  {x:\Int\vdash x<10:\Bool}}
{\emptyset\vdash \lambda(x:\Int).\ x<10:\Int\Arrow\Bool}
\]
\begin{itemize}
  \item Add the assumption $x:\Int$ while checking the body.
  \item Both sides of $<$ are integers, so the body is boolean.
  \item The whole function therefore has type $\Int\Arrow\Bool$.
\end{itemize}
\end{frame}

\begin{frame}{Work one through: call the function}
Let $\mathit{small}=\lambda(x:\Int).\ x<10$. What is $\mathit{small}\ 4$?
\[
\frac{\emptyset\vdash\mathit{small}:\Int\Arrow\Bool\qquad
      \emptyset\vdash 4:\Int}
     {\emptyset\vdash\mathit{small}\ 4:\Bool}
\]
Now try $\mathit{small}\ \mathsf{false}$:
\begin{alertblock}{Mismatch}
The function expects an $\Int$, but the argument has type $\Bool$. Reject the call before execution.
\end{alertblock}
\end{frame}

\begin{frame}{How a checker can do the work}
For this language, checking follows the syntax tree recursively:
\begin{enumerate}
  \item A literal has its built-in type; a variable uses the environment.
  \item For $e_1+e_2$, check both children are $\Int$; return $\Int$.
  \item For a conditional, check a $\Bool$ condition and matching branch types.
  \item For a function, extend the environment with its parameter type.
  \item For a call, check the function type and the argument type agree.
\end{enumerate}
\begin{block}{A useful mental model}
The rules are a recipe for a recursive program: inspect the outermost expression, check its parts, and combine their types.
\end{block}
\end{frame}

\begin{frame}{When the recipe says “no”}
\begin{columns}[T]
\column{.49\textwidth}
\begin{alertblock}{Bad condition}
\texttt{if 3 then ... else ...}\
Expected $\Bool$; got $\Int$.
\end{alertblock}
\begin{alertblock}{Branches disagree}
\texttt{if ready then 0 else false}\
One branch is $\Int$, the other $\Bool$.
\end{alertblock}
\column{.49\textwidth}
\begin{alertblock}{Bad operator input}
\texttt{true + 1}\
$+$ requires two integers.
\end{alertblock}
\begin{alertblock}{Bad function argument}
\texttt{(fun (n:Int) => n+1) false}\
Expected $\Int$; got $\Bool$.
\end{alertblock}
\end{columns}
\medskip
Good diagnostics point to the expression and explain the expected and actual types.
\end{frame}

\begin{frame}{A reassuring promise: progress (informally)}
For a well-typed program in a carefully designed language:
\begin{itemize}
  \item Evaluation does not get stuck trying to add a boolean to an integer.
  \item A function call does not apply a boolean as though it were a function.
\end{itemize}
\begin{block}{This is the type-safety idea}
Types rule out a useful class of runtime failures. The formal result is often called \emph{type soundness}; its details depend on the language and its operations.
\end{block}
It does not rule out every runtime failure: division by zero, bad input, or nontermination may still be possible.
\end{frame}

\begin{frame}{Try it together}
Assume $f:\Int\Arrow\Bool$ and $n:\Int$. Find the type, or the first reason it fails:
\begin{enumerate}
  \item $f\ (n+1)$
  \item $\mathsf{if}\ f\ n\ \mathsf{then}\ n+2\ \mathsf{else}\ 0$
  \item $\mathsf{if}\ n\ \mathsf{then}\ \mathsf{true}\ \mathsf{else}\ \mathsf{false}$
\end{enumerate}
\pause
\begin{exampleblock}{Answers}
  \begin{enumerate}
    \item $\Bool$.
    \item $\Int$.
    \item Rejected: condition $n$ has type $\Int$, but needs $\Bool$.
  \end{enumerate}
\end{exampleblock}
\end{frame}

\begin{frame}{Takeaways}
\begin{itemize}
  \item A type describes how an expression may be used.
  \item A typing judgment combines an expression, its environment, and its type.
  \item Each syntax form gets a small rule; checking is recursive.
  \item Function types connect the argument you provide to the result you get.
  \item Simple types catch real mistakes, while leaving plenty of interesting bugs for later.
\end{itemize}
\bigskip
\centering\Large Questions?
\end{frame}

\end{document}
